%================================================== Описание и работа ==================================================
\color{unidarkgreen}\section[Описание и работа]{Описание и работа}\color{black}


%================================================== Назначение ==================================================
\color{unidarkgreen}{\subsection{Назначение}}\color{black}

\begin{enumerate}[label=\arabic{section}.\arabic{subsection}.\arabic*, labelsep=4pt, leftmargin=0pt, itemindent=2cm, align=left]

\item
Шкаф «ШЭТ 251.01-0-ЮИРЗ» предназначен для выполнения функций автоматики управления выключателем (АУВ) и устройства резервирования отказа выключателя (УРОВ) на электрических станциях и подстанциях напряжением 110-220 кВ архитектуры I типа при новом строительстве или реконструкции энергообъктов.

\item
Шкаф предназначен для устанвоки в помещениях релейных щитов и общеподстанционных пунктов управления.

\item
Исполнение шкафа может быть выполнено с переменным, выпрямленным переменным или постоянным оперативным током в соответствии с указанием в карте заказа {\color{red}(см. приложение КАРТА ЗАКАЗА)}.

\item
Функциональное назначение шкафа отражается в структуре его условного обозначения в соответствии с СТО 56947007-33.040.20.285-2019 ПАО «ФСК ЕЭС».

\item
Исполнение шкафа выполнено с установкой в нем одного интеллектуального электронного устройства (ИЭУ) серии «ЮНИТ-М500», имеющего код заказа ЮНИТ-М514-ЛВ-Р102-К103-К103-К103-В101-В101-В101-М179.0700-С01.03 (см. приложение \ref{app:card}).

\item
Функциональное назначение ИЭУ:
\medskip
\begin{itemize}
\item контроль синхронизма и напряжений;
\item автоматическое повторное включение (АПВ);
\item полуавтоматическое включение (ПАВ);
\item контроль силового выключателя (КСВ);
\item функциональный блок управления выключателем;
\item Функциональный блок фиксации положения выключателей (ФПВ);
\item функциональный блок фиксации состояния ЛЭП (ФСЛ) 
\item устройство резервирования отказа выключателя (УРОВ).
\end{itemize}
\medskip

\item
Дополнительно ИЭУ выполняет функции:
\medskip
\begin{itemize}
\item местной предупредительной и аварийной сигнализации;
\item регистрации аварийных событий (РАС);
\item самодиагностики.
\end{itemize}
\medskip

\end{enumerate}


%================================================== Технические характеристики ==================================================
\color{unidarkgreen}\subsection{Технические характеристики}\color{black}
\subsubsection{Основные параметры}

\begin{enumerate}[label=\arabic{section}.\arabic{subsection}.\arabic{subsubsection}.\arabic*,  align=left, labelsep=4pt, leftmargin=0pt, itemindent=57pt]

\item
Номинальные параметры шкафа указаны в таблице \ref{sec1:tbl1}

\setlength{\extrarowheight}{0.06cm} % высота ячеек в таблице
\small
\begin{longtable}
{|>{\centering\arraybackslash}m{12.9cm}|>{\centering\arraybackslash}m{4cm}|}
\caption{Основные номинальные параметры шкафа\hfill\vspace{-0.5\baselineskip}}\label{sec1:tbl1}\\ 
\hhline{|-|-|}
\rowcolor{gray!30}
Параметр & Значение \\
\hhline{|-|-|}
\endfirsthead

%\caption{\emph{Продолжение\hfill\vspace{-0.5\baselineskip}}} \\ % Отступ обозначения таблицы от таблицы и выравнивание надписи слева
\caption*{\hspace{3pt}\emph{Продолжение таблицы \ref{sec1:tbl1}\hfill\vspace{-0.5\baselineskip}}} \\ % сделано по ГОСТ 2.105 п.6.8.7
\hhline{|-|-|}
\rowcolor{gray!30}
Параметр & Значение \\ 
\hhline{|-|-|}
\endhead
\hhline{|-|-|}

\endfoot
\endlastfoot

\raggedright Номинальный переменный ток, А  & \centering\arraybackslash 1 или 5 \\  \hhline{|-|-|}
\raggedright Номинальное линейное напряжение переменного тока, В  & \centering\arraybackslash 100 \\  \hhline{|-|-|}
\raggedright Номинальная частота, Гц  & \centering\arraybackslash 50 \\  \hhline{|-|-|}
\raggedright Номинальное напряжение оперативного (постоянного / переменного / выпрямленного переменного) тока, В & \centering\arraybackslash 110 или 220 \\  \hhline{|-|-|}

\end{longtable}
\normalsize
\setlength{\extrarowheight}{0.0cm} % возврат параметра высота ячеек в таблице

\end{enumerate}

\subsubsection{Условия эксплуатации}

\begin{enumerate}[label=\arabic{section}.\arabic{subsection}.\arabic{subsubsection}.\arabic*,  align=left, labelsep=4pt, leftmargin=0pt, itemindent=57pt]

\item
Климатическое исполнение шкафа и категория его размещения по ГОСТ 15150-69 – УХЛ3.1(УХЛ4).

\item
Условия эксплуатации в соответствии с ГОСТ IEC 61439-1:
\begin{itemize}
\item температура окружающей среды должна быть не более 40°С, а средняя температура за 24 ч -- не более 35°С;
\item минимальное значение температуры окружающей среды -- минус 5°С;
\item относительная влажность воздуха не должна превышать 50\% при максимальной температуре 40°С. При более низких температурах допускается более высокая относительная влажность, например 90\% при 20°С;
\item высота установки над уровнем моря не должна превышать 2000 м;
\item степень загрязнения окружающей среды 1 -- загрязнение отсутствует или имеется только сухое непроводящее загрязнение. Загрязнение незначительное;
\item место установки устройства должно быть защищено от попадания брызг воды, масел, эмульсий, а также от прямого воздействия солнечной радиации.
\end{itemize}

\item
Рабочее положение шкафа в пространстве -- вертикальное с отклонением от рабочего положения до 5° в любую сторону.

\item
В части воздействия факторов внешней среды шкаф удовлетворяет требованиям группы механического исполнения М40 по ГОСТ 30631-99. Уровень вибрационных нагрузок от 0,5 до 100 Гц с ускорением 0,5 g. Шкаф выдерживает многократные ударные нагрузки длительностью от 2 до 20 мс с максимальным ускорением 3 g.

\item
Шкаф сейсмостоек при воздействии землетрясений интенсивностью до 9 баллов включительно по шкале MSK-64 при уровне установки над нулевой отметкой до 10 м по ГОСТ 17516.1-90.

\item
Степень защиты, обеспечиваемая оболочкой шкафа, по ГОСТ 14254-2015 от доступа к опасным частям, попадания внешних твердых предметов и (или) воды соответствует IP54.

\end{enumerate}

\subsubsection{Сопротивление и электрическая прочность изоляции}

\begin{enumerate}[label=\arabic{section}.\arabic{subsection}.\arabic{subsubsection}.\arabic*,  align=left, labelsep=4pt, leftmargin=0pt, itemindent=57pt]

\item
Сопротивление изоляции электрически независимых цепей шкафа, измеренное в холодном состоянии при температуре окружающего воздуха 20±5°С и относительной влажности до 80\%, должно быть не менее 100 МОм при напряжении постоянного тока 500 В.

\item
В состоянии поставки электрическая изоляция между всеми независимыми цепями шкафа (кроме портов последовательной передачи данных терминалов) относительно корпуса и всех независимых цепей между собой выдерживает без пробоя и перекрытия испытательное напряжение, приведенное в таблице \ref{sec1:tbl2}.

\setlength{\extrarowheight}{0.06cm} % высота ячеек в таблице
\small
\begin{longtable}
{|>{\centering\arraybackslash}m{12.9cm}|>{\centering\arraybackslash}m{4cm}|}
\caption{Требования к испытанию изоляции\hfill\vspace{-0.5\baselineskip}}\label{sec1:tbl2}\\ 
\hhline{|-|-|}
\rowcolor{gray!30}
Наименование показателя & Значение \\
\hhline{|-|-|}
\endfirsthead

\caption*{\hspace{3pt}\emph{Продолжение таблицы \ref{sec1:tbl2}\hfill\vspace{-0.5\baselineskip}}} \\ % сделано по ГОСТ 2.105 п.6.8.7
\hhline{|-|-|}
\rowcolor{gray!30}
Наименование показателя & Значение \\ 
\hhline{|-|-|}
\endhead
\hhline{|-|-|}

\endfoot
\endlastfoot

\raggedright Электрическая прочность изоляции цепей с рабочим напряжением не более 60 В   & \centering\arraybackslash 500 В действующего значения, 50 Гц, 1 мин \\  \hhline{|-|-|}
\raggedright Электрическая прочность изоляции цепей с рабочим напряжением более 60 В   & \centering\arraybackslash 2000 В действующего значения, 50 Гц, 1 мин \\  \hhline{|-|-|}

\end{longtable}
\normalsize
\setlength{\extrarowheight}{0.0cm} % возврат параметра высота ячеек в таблице

\item
При повторных испытаниях шкафа испытательное напряжение не должно превышать 85\% от вышеуказанных значений.

\end{enumerate}

\subsubsection{Цепи оперативного питания}

\begin{enumerate}[label=\arabic{section}.\arabic{subsection}.\arabic{subsubsection}.\arabic*,  align=left, labelsep=4pt, leftmargin=0pt, itemindent=57pt]

\item
Питание шкафа осуществляется от цепей оперативного тока.
Параметры цепей питания ИЭУ, входящего в состав шкафа, приведены в руководстве по эксплуатации {\color{red}«УСТРОЙСТВО ЗАЩИТЫ И АВТОМАТИКИ СЕРИИ «ЮНИТ-М500» РУКОВОДСТВО ПО ЭКСПЛУАТАЦИИ ОБЩЕЕ»}.

\item
Характеристики дискретных входов и выходов ИЭУ, входящего в состав шкафа, приведены в руководстве по эксплуатации {\color{red}«УСТРОЙСТВО ЗАЩИТЫ И АВТОМАТИКИ СЕРИИ «ЮНИТ-М500» РУКОВОДСТВО ПО ЭКСПЛУАТАЦИИ ОБЩЕЕ»}.

\item
Шкаф правильно функционирует при изменении напряжения оперативного тока в диапазоне от 0,8 до 1,1 номинального значения.

\item
Контакты выходных реле шкафа не замыкаются ложно при подаче и снятии напряжения оперативного тока с перерывом любой длительности.

\end{enumerate}

\subsubsection{Аналоговые цепи}

\begin{enumerate}[label=\arabic{section}.\arabic{subsection}.\arabic{subsubsection}.\arabic*,  align=left, labelsep=4pt, leftmargin=0pt, itemindent=57pt]

\item
Характеристики аналоговых цепей ИЭУ, входящего в состав шкафа, приведены в руководстве по эксплуатации {\color{red}«УСТРОЙСТВО ЗАЩИТЫ И АВТОМАТИКИ СЕРИИ «ЮНИТ-М500» РУКОВОДСТВО ПО ЭКСПЛУАТАЦИИ ОБЩЕЕ»}.

\end{enumerate}

\subsubsection{Электромагнитная совместимость}

\begin{enumerate}[label=\arabic{section}.\arabic{subsection}.\arabic{subsubsection}.\arabic*,  align=left, labelsep=4pt, leftmargin=0pt, itemindent=57pt]

\item
Требования по электромагнитной совместимости согласно приложению J ГОСТ IEC 61439-1-2024. Требования обеспечиваются входящими в шкаф компонентами. Устойчивость терминала к воздействию помех приведена в руководстве по эксплуатации на терминал {\color{red}«УСТРОЙСТВО ЗАЩИТЫ И АВТОМАТИКИ СЕРИИ «ЮНИТ-М500» РУКОВОДСТВО ПО ЭКСПЛУАТАЦИИ ОБЩЕЕ»}.

\end{enumerate}

\subsubsection{Надежность}

\begin{enumerate}[label=\arabic{section}.\arabic{subsection}.\arabic{subsubsection}.\arabic*,  align=left, labelsep=4pt, leftmargin=0pt, itemindent=57pt]

\item
Средняя наработка на отказ шкафа составляет не менее 125 000 ч.

\item
Среднее время восстановления работоспособного состояния ИЭУ при наличии полного комплекта запасных блоков составляет не более 2 ч с учетом времени нахождения неисправности.

\item
Средний срок службы шкафа - не менее 25 лет при условии проведения требуемых технических мероприятий по обслуживанию с заменой, при необходимости, материалов и комплектующих, имеющих меньший срок службы.

\item
Средний срок сохраняемости шкафа в упаковке поставщика составляет три года.

\item
Периодичность планового технического обслуживания определяется характеристиками ИЭУ и комплектующих в ходящих в состав шкафа, а также моделью организации работ в эксплуатирующей организации:
\begin{itemize}
\item по графику;
\item по состоянию, с автоматизированным дистанционным мониторингом или без него.
\end{itemize}

\item
Минимальный срок периодического технического обслуживания шкафов должен составлять не менее 3 лет.

\end{enumerate}

\subsubsection{Защитное заземление}

\begin{enumerate}[label=\arabic{section}.\arabic{subsection}.\arabic{subsubsection}.\arabic*,  align=left, labelsep=4pt, leftmargin=0pt, itemindent=57pt]

\item
Защитное заземление шкафа осуществляется через отдельный заземляющий проводник подключенный к специальному элементу заземления.

\item
Шкаф имеет элементы для заземления (шина, болты, винты).
Диаметр болтов (винтов, шпилек) для заземления, размеры контактной площадки, к которым присоединяются защитные провода заземления, должны соответствовать ГОСТ 12.2.007.0.
Элементы заземления размещаются внутри шкафа.

\item
Все металлические части устройств в составе шкафа и металлоконструкции, которые могут оказаться под напряжением вследствие нарушения изоляции, заземляются.

\item
В соответствии с ГОСТ 12.2.007.0 обеспечивается непрерывность цепи защитного заземления.
Значение сопротивления между заземляющим болтом для заземления шкафа и каждой доступной прикосновению металлической нетоковедущей частью оборудования не должно превышать 0.1 Ом (заземляющая цепь должна быть не прерывной).

\item
Проводники защитного заземления отличаются от других цепей и имеют двухцветную изоляцию желто-зеленого цвета по всей своей длине.
Главный элемент заземления шкафа имеет знак «заземления».

\end{enumerate}


%================================================== Состав шкафа ==================================================
\color{unidarkgreen}\subsection{Состав шкафа}\color{black}
\subsubsection{Конструктивное исполнение}

\begin{enumerate}[label=\arabic{section}.\arabic{subsection}.\arabic{subsubsection}.\arabic*,  align=left, labelsep=4pt, leftmargin=0pt, itemindent=57pt]

\item
Габаритные и установочные размеры шкафа представлены в {\color{red}приложении \ref{app:size}}.

\item
Шкаф представляет собой защищённое низковольтное комплектное устройство (НКУ).
Шкаф представляет собой металлоконструкцию из специализированного профиля с двухсторонним обслуживанием.

\item
Габариты шкафа 800х600х2200 мм (ШхГхВ, включая цоколь 200 мм).

\item
Несущей конструкцией изделия является цельносварная или сборная рама.
Для крепления изделия применяется монтажная конструкция в виде цоколя.
Изделия крепятся к цоколю болтовыми соединениями.

\item
Для крепления оборудования внутри шкафа предусматриваются внутренние поворотные рамы, двери или монтажные панели.

\item
Оболочка шкафа предусматривает двери: лицевая (одностворчатая - сплошная металлическая, со смотровым окном или сплошная обзорная) и задняя (двухстворчатая или одностворчатая - глухая).

\item
Двери обеспечивают свободное открывание и закрывание на угол не менее 110°, задняя дверь оборудована упорами для фиксации двери в открытом положении.
Двери шкафов снабжены замками, предотвращающими несанкционированный доступ.

\color{red}

\item
На передней двери в исполнении - сплошная металлическая, со смотровым окном, располагаются: 
\begin{itemize}
\item лампы сигнализации;
\item оперативные переключатели или блоки электронных ключей;
\item кнопки управления (квитирования, деблокировки и т.п.;
\item указательные реле.
\end{itemize}

\item
На монтажной плите шкафа в зависимости от типоисполнения шкафа и исполнения передней двери могут располагаться:
\begin{itemize}
\item терминал(ы);
\item ВЧ-приемопередатчик;
\item блоки испытательные (SG), через которые подключаются входные аналоговые цепи шкафа от трансформаторов тока и напряжения;
\item оперативные переключатели или блоки электронных ключей;
\item кнопки управления;
\item лампы сигнализации;
\item указательные реле;
\item автоматы питания (SF), через которые подключаются независимые цепи оперативного питания «±ЕС».
\end{itemize}

\color{black}

\item
Подвод кабельных связей к шкафу осуществляется снизу через дно шкафа.
Кабельные связи подключаются к клеммным рядам, расположенные на провой и/или левой боковине.
Доступ к клеммным рядам, устройствам и монтажу шкафа обеспечивается при открытие задней двери, для двухстороннего облуживания и поворотной рамы, для одностороннего обслуживания. 

Монтаж аппаратов шкафа между собой выполнен медными проводами.
Номинальное сечение проводов не менее 2,5~мм$^2$ для токовых цепей, не менее 0,75~мм$^2$ - для остальных цепей.
Допускается отклонение от указанных требований при условии обеспечения выполнения требований к термической стойкости и механической прочности.

\item
Для цепей тока допускается подключение одного проводника сечением не более 10~мм$^2$ или двух проводников сечением не более 2,5~мм$^2$.
Для остальных цепей допускается подключение одного проводника сечением не более 6~мм$^2$ или двух проводников сечением не более 1,5~мм$^2$.

\end{enumerate}

\subsubsection{Схемы шкафа}
\begin{enumerate}[label=\arabic{section}.\arabic{subsection}.\arabic{subsubsection}.\arabic*,  align=left, labelsep=4pt, leftmargin=0pt, itemindent=57pt]
\color{red}
\item Схемы шкафа представлены в приложении Г.;
\item Схема электрическая принципиальная ЮТКБ.ХХ-ХХХХ.ХХХ Э3;
\item Схема электрическая монтажная ЮТКБ. ХХ-ХХХХ.ХХХ Э4;
\item Схема рядов зажимов ЮТКБ. ХХ-ХХХХ.ХХХ Э4.1.
\color{black}
\end{enumerate}


%================================================== Устройство и работа ==================================================
\color{unidarkgreen}\subsection{Устройство и работа}\color{black}
\begin{enumerate}[label=\arabic{section}.\arabic{subsection}.\arabic*,  align=left, labelsep=4pt, leftmargin=0pt, itemindent=57pt]
\color{red}
\item
Принцип работы функций защит и автоматики ИЭУ, входящих в состав шкафа, описаны в руководстве по эксплуатации {\color{red} «Микропроцессорное устройство. Автоматика управления выключателем 110-220 кВ, УРОВ. Руководство по эксплуатации.»}.
\color{black}
\end{enumerate}


%================================================== Средства измерений, инструмент и принадлежности ==================================================
\color{unidarkgreen}\subsection{Средства измерений, инструмент и принадлежности}\color{black}
\begin{enumerate}[label=\arabic{section}.\arabic{subsection}.\arabic*,  align=left, labelsep=4pt, leftmargin=0pt, itemindent=57pt]

\item
Перечень оборудования и средств измерений, необходимых для проведения эксплуатационных проверок шкафа, приведён в приложении \ref{app:unit}.

\end{enumerate}


%================================================== Маркировка и пломбирование ==================================================
\color{unidarkgreen}\subsection{Маркировка и пломбирование}\color{black}
Шкаф имеет маркировку в соответствии с конструкторской документацией. Маркировка выполнена в соответствии с ГОСТ 18620-86 способом, обеспечивающим её чёткость и сохраняемость.

\par
Информационная табличка (шильдик) размещается в правом верхнем углу на передней двери шкафа и дублируется на монтажной панели шкафа с лицевой стороны. На табличке (шильдике) в дополнение к текстовой информации наносится QR-код содержащий:
\begin{itemize}
\item наименование шкафа;
\item шифр шкафа;
\item основные функции МП ИЭУ шкафа;
\item номинальный вторичный переменный ток;
\item номинальная частота;
\item номинальное переменное напряжение;
\item напряжение оперативного постоянного тока;
\item дата (месяц, год) выпуска шкафа в формате ММ.ГГГГ.
\end{itemize}

\par
Все элементы шкафа имеют обозначение, соответствующее схеме шкафа и состоящее из буквенного обозначения и порядкового номера, проставленного после буквенного обозначения (например, SG1).

\par
Транспортная маркировка тары выполнена в соответствие с ГОСТ 14192-96.
Манипуляционные знаки наносят в верхней части двух вертикальных смежных сторон.


%================================================== Упаковка ==================================================
\color{unidarkgreen}\subsection{Упаковка}\color{black}
Упаковка шкафа проводится по конструкторской документации предприятия-изготовителя для условий хранения, транспортирования и допустимых сроков сохраняемости.

\par
В соответствии с ГОСТ 23216 для изделий применяется транспортная тара двух видов:
\begin{itemize}
\item ТЭ: ящики дощатые;
\item ТФ: ящики фанерные и из древесно-волокнистых плит.
\end{itemize}
